\sectionkicker{Source-backed technical notes}
\subsection*{巨大モデルを作り直さずに変える――Prefix、Prompt Tuning、LoRA: Technical Notes}
本欄は記事本文で圧縮した一次資料上の情報を、比較・再検証しやすい形で整理したものである。時系列、確認済みの主張、留意点、一次資料URLを掲載する。
\smallskip
\noindent{\small\textit{Technical Notesでは、一次資料から確認した公開・提供・研究上の事実を記載する。数値・能力評価や提供元・著者による比較表現は、独立再現済みの結果ではなく帰属claimとして扱う。}}\par
\medskip
\subsection*{Theme at a glance}
\addcontentsline{toc}{subsection}{Theme at a glance}
\begin{center}
\footnotesize
\begin{tabularx}{\linewidth}{@{}p{0.20\linewidth}p{0.10\linewidth}p{0.14\linewidth}X@{}}
\toprule
資料 & 位置づけ & 種別 & 時系列 \\
\midrule
Prefix-Tuning & 主要資料 & 論文 & 2021-01-01 \\
Prompt Tuning & 主要資料 & 論文 & 2021-04-18 \\
LoRA & 主要資料 & 論文 & 2021-06-17 \\
\bottomrule
\end{tabularx}
\end{center}
\normalsize

\begin{technicalnote}{Prefix-Tuning}{主要資料}
\begingroup
% reader-facing Technical Notes break policy
\widowpenalty=10000
\clubpenalty=10000
\displaywidowpenalty=10000
\begin{tabularx}{\linewidth}{@{}>{\bfseries}p{0.22\linewidth}X@{}}
組織 & Stanford / authors \\
種別 & 論文 \\
時系列 & 2021-01-01 \\
\end{tabularx}
\smallskip
\begin{minipage}{\linewidth}
% reader-facing Technical Notes event-only fact/source tail group
{\bfseries 一次資料から整理したtechnical points}
\begin{itemize}[leftmargin=1.5em,itemsep=0.35em]
\item \textbf{一次情報で確認できる事実}: Stanford / authorsの選定済み一次資料では、Prefix-Tuningはpretrained language model本体のparameterをfreezeしたまま、task-specificな小さなcontinuous vector（prefix）だけを最適化し、後続tokenがそのprefixをvirtual tokenのように参照するparameter-efficient adaptationとして提案された。 論文はGPT-2のtable-to-text generationとBARTのsummarizationを評価対象とし、著者らは全parameterの0.1\%のみを学習する設定を含めてfull-data・low-data・unseen-topicへの一般化を比較した。
\end{itemize}
\begin{minipage}{\linewidth}
% reader-facing Technical Notes generic boundary/source tail group
\begin{samepage}
{\bfseries 一次資料}
\begingroup\sloppy
\Urlmuskip=0mu plus 2mu\relax
\begin{itemize}[leftmargin=1.5em,itemsep=0.25em]
\item {\scriptsize\url{http://arxiv.org/abs/2101.00190v1}}
\end{itemize}
\endgroup
\end{samepage}
\end{minipage}
\end{minipage}
\endgroup
\end{technicalnote}
\medskip

\begin{technicalnote}{Prompt Tuning}{主要資料}
\begingroup
% reader-facing Technical Notes break policy
\widowpenalty=10000
\clubpenalty=10000
\displaywidowpenalty=10000
\begin{tabularx}{\linewidth}{@{}>{\bfseries}p{0.22\linewidth}X@{}}
組織 & Google / authors \\
種別 & 論文 \\
時系列 & 2021-04-18 \\
\end{tabularx}
\smallskip
\begin{minipage}{\linewidth}
% reader-facing Technical Notes event-only fact/source tail group
{\bfseries 一次資料から整理したtechnical points}
\begin{itemize}[leftmargin=1.5em,itemsep=0.35em]
\item \textbf{一次情報で確認できる事実}: Google / authorsの選定済み一次資料では、Prompt Tuningはlanguage model本体をfreezeしたまま、backpropagationで学習するsoft promptによってdownstream taskへ条件付けするparameter-efficient adaptationとして提案された。 著者らはT5でmodel sizeを変えたablationを行い、modelがbillion-parameter規模を超えるにつれてprompt tuningがfull model tuningとの差を縮めること、およびdomain transferでのrobustnessを評価した。
\end{itemize}
\begin{minipage}{\linewidth}
% reader-facing Technical Notes generic boundary/source tail group
\begin{samepage}
{\bfseries 一次資料}
\begingroup\sloppy
\Urlmuskip=0mu plus 2mu\relax
\begin{itemize}[leftmargin=1.5em,itemsep=0.25em]
\item {\scriptsize\url{http://arxiv.org/abs/2104.08691v2}}
\end{itemize}
\endgroup
\end{samepage}
\end{minipage}
\end{minipage}
\endgroup
\end{technicalnote}
\medskip

\begin{technicalnote}{LoRA}{主要資料}
\begingroup
% reader-facing Technical Notes break policy
\widowpenalty=10000
\clubpenalty=10000
\displaywidowpenalty=10000
\begin{tabularx}{\linewidth}{@{}>{\bfseries}p{0.22\linewidth}X@{}}
組織 & Microsoft / authors \\
種別 & 論文 \\
時系列 & 2021-06-17 \\
\end{tabularx}
\smallskip
\begin{minipage}{\linewidth}
% reader-facing Technical Notes event-only fact/source tail group
{\bfseries 一次資料から整理したtechnical points}
\begin{itemize}[leftmargin=1.5em,itemsep=0.35em]
\item \textbf{一次情報で確認できる事実}: Microsoft / authorsの選定済み一次資料では、LoRAはpretrained weightをfreezeし、各layerのweight updateを低rank行列の積としてparameter化して、その追加parameterだけを学習するadaptation手法として提案された。 論文はGPT-2・GPT-3・RoBERTa・DeBERTaでfull fine-tuningやadapter系手法と比較し、trainable parameter数・memory・inference latencyを含む運用上の差も評価した。
\end{itemize}
\begin{minipage}{\linewidth}
% reader-facing Technical Notes generic boundary/source tail group
\begin{samepage}
{\bfseries 一次資料}
\begingroup\sloppy
\Urlmuskip=0mu plus 2mu\relax
\begin{itemize}[leftmargin=1.5em,itemsep=0.25em]
\item {\scriptsize\url{http://arxiv.org/abs/2106.09685v2}}
\end{itemize}
\endgroup
\end{samepage}
\end{minipage}
\end{minipage}
\endgroup
\end{technicalnote}
\medskip
